% 03-lexer.tex — Chapter 3: Tokens and Lexer
\chapter{Tokens and Lexer}
\label{chap:03-lexer}
\index{lexer}
\index{tokeniser}

\pnum
The Lain lexer converts a NUL-terminated source buffer into a flat stream of
\cname{Token} values.  It is implemented in \srcref{src/lexer.h} and uses
\srcref{src/token.h} for the token-kind enumeration and keyword-matching logic.

\pnum
The lexer is zero-copy: every token holds a \texttt{const char *} pointer into
the original source buffer together with a byte length.  No new strings are
allocated during lexing.  The source buffer lives in \cname{file\_arena} for
the lifetime of the compiler run.

% -------------------------------------------------------------------------
\section{Token representation}
\label{sec:lex-token}
\index{Token struct}
% -------------------------------------------------------------------------

\pnum
A \cname{Token} is defined in \srcref{src/token.h} as:

\begin{ccode}
typedef struct {
    TokenKind   kind;
    const char *start;
    isize       length;
} Token;
\end{ccode}

\begin{structdoc}{Token}{src/token.h:97}
\cname{kind}   & \cname{TokenKind} & Discriminant identifying the token category. \\
\cname{start}  & \cname{const char *} & Pointer to the first byte of the lexeme in the source buffer. \\
\cname{length} & \cname{isize}        & Byte length of the lexeme. \\
\end{structdoc}

\pnum
String tokens (\cname{TOKEN\_STRING\_LITERAL}) are an exception: \cname{start}
points to the first byte \emph{after} the opening double-quote and \cname{length}
excludes both delimiters.  This makes string content immediately available
without any stripping step in the parser.

% -------------------------------------------------------------------------
\section{Token kinds}
\label{sec:lex-kinds}
\index{TokenKind}
% -------------------------------------------------------------------------

\pnum
The \cname{TokenKind} enumeration (\srcref{src/token.h:6}) contains 95~members
grouped as follows:

\begin{center}
\begin{tabular}{ll}
\toprule
\textbf{Group} & \textbf{Representative members} \\
\midrule
Control      & \cname{TOKEN\_EOF}, \cname{TOKEN\_NEWLINE}, \cname{TOKEN\_INVALID} \\
Literals     & \cname{TOKEN\_NUMBER}, \cname{TOKEN\_FLOAT\_LITERAL},
               \cname{TOKEN\_STRING\_LITERAL}, \cname{TOKEN\_CHAR\_LITERAL} \\
Identifiers  & \cname{TOKEN\_IDENTIFIER} \\
Delimiters   & \cname{TOKEN\_L\_PAREN}, \cname{TOKEN\_R\_BRACE},
               \cname{TOKEN\_L\_BRACKET}, \ldots \\
Operators    & \cname{TOKEN\_PLUS}, \cname{TOKEN\_MINUS},
               \cname{TOKEN\_ASTERISK}, \ldots \\
Compound ops & \cname{TOKEN\_PLUS\_EQUAL}, \cname{TOKEN\_EQUAL\_EQUAL},
               \cname{TOKEN\_SHIFT\_LEFT}, \ldots \\
Punctuation  & \cname{TOKEN\_DOT}, \cname{TOKEN\_DOT\_DOT},
               \cname{TOKEN\_DOT\_DOT\_EQUAL}, \cname{TOKEN\_COLON}, \ldots \\
Comments     & \cname{TOKEN\_LINE\_COMMENT}, \cname{TOKEN\_MULTILINE\_COMMENT} \\
Keywords     & \cname{TOKEN\_KEYWORD\_IF} through \cname{TOKEN\_KEYWORD\_DECREASING} \\
\bottomrule
\end{tabular}
\end{center}

% -------------------------------------------------------------------------
\section{Lexer state machine}
\label{sec:lex-fsm}
\index{lexer!FSM}
\index{LexerState}
% -------------------------------------------------------------------------

\pnum
The \cname{Lexer} struct holds two pointers:

\begin{structdoc}{Lexer}{src/lexer.h:9}
\cname{text}    & \cname{const char *} & Pointer to the first byte of the source buffer (immutable). \\
\cname{current} & \cname{const char *} & Pointer to the next byte to be consumed. \\
\end{structdoc}

\pnum
The core function \cname{lexer\_next(Lexer *lexer)} (\srcref{src/lexer.h:51})
implements a single-function FSM.  On entry, \cname{token.start} is set to
\cname{current}.  The FSM loops over bytes, transitions states via
\cname{LexerState}, and fires one of two macros:

\begin{ccode}
#define SWITCH_STATE(_state) state = _state; break
#define RETURN_TOKEN(_kind)  token.kind = _kind; \
                             token.length = lexer->current - token.start; \
                             return token
\end{ccode}

\cname{SWITCH\_STATE} sets the next state and breaks the inner switch so the
outer loop reads the next byte.  \cname{RETURN\_TOKEN} computes the token length
and returns.

\pnum
The 22-member \cname{LexerState} enum:

\begin{center}
\begin{tabular}{ll}
\toprule
\textbf{State} & \textbf{Entered after} \\
\midrule
\cname{STATE\_START}                & Start of every token \\
\cname{STATE\_IDENTIFIER}           & Any letter or \texttt{\_} \\
\cname{STATE\_NUMBER}               & Any decimal digit \\
\cname{STATE\_SINGLE\_QUOTE}        & \texttt{'} \\
\cname{STATE\_DOUBLE\_QUOTE}        & \texttt{"} \\
\cname{STATE\_LINE\_COMMENT}        & \texttt{//} \\
\cname{STATE\_EQUAL}                & \texttt{=} \\
\cname{STATE\_ANGLE\_BRACKET\_LEFT} & \texttt{<} (produces \texttt{<}, \texttt{<=}, \texttt{<<}) \\
\cname{STATE\_ANGLE\_BRACKET\_RIGHT}& \texttt{>} (produces \texttt{>}, \texttt{>=}, \texttt{>>}) \\
\cname{STATE\_DOT}                  & \texttt{.} \\
\cname{STATE\_DOT\_DOT}             & \texttt{..} \\
\cname{STATE\_ASTERISK}             & \texttt{*} \\
\cname{STATE\_SLASH}                & \texttt{/} (dispatches to \texttt{/}, \texttt{/=}, \texttt{//}, \texttt{/*\!}) \\
\cname{STATE\_PERCENT}              & \texttt{\%} \\
\cname{STATE\_PLUS}                 & \texttt{+} \\
\cname{STATE\_MINUS}                & \texttt{-} \\
\cname{STATE\_BANG}                 & \texttt{!} \\
\cname{STATE\_AMPERSAND}            & \texttt{\&} \\
\cname{STATE\_PIPE}                 & \texttt{|} \\
\cname{STATE\_CARET}                & \texttt{\^{}} \\
\bottomrule
\end{tabular}
\end{center}

% -------------------------------------------------------------------------
\section{Special lexing cases}
\label{sec:lex-special}
% -------------------------------------------------------------------------

\subsection{Whitespace and newlines}
\label{sec:lex-whitespace}

\pnum
Horizontal whitespace (\texttt{space} and \texttt{tab}) in \cname{STATE\_START}
advances \cname{token.start} silently, skipping the character without emitting a
token.  Vertical whitespace (\texttt{\\n} and \texttt{\\r}) returns
\cname{TOKEN\_NEWLINE}.  The parser uses newlines as statement terminators in
contexts where a semicolon would otherwise be required.

\subsection{Numeric literals}
\label{sec:lex-numbers}
\index{numeric literals!lexing}

\pnum
Numeric scanning (\cname{STATE\_NUMBER}) handles four bases:
\begin{itemize}
\item Decimal: sequences of \texttt{0}--\texttt{9} and underscore separators.
\item Hexadecimal: prefix \texttt{0x}/\texttt{0X} followed by hex digits and underscores.
\item Binary: prefix \texttt{0b}/\texttt{0B} followed by \texttt{0}/\texttt{1} and underscores.
\item Octal: prefix \texttt{0o}/\texttt{0O} followed by octal digits and underscores.
\end{itemize}

\pnum
A decimal point followed by digits produces \cname{TOKEN\_FLOAT\_LITERAL}.
If the decimal point is immediately followed by another dot (i.e.\ \texttt{..}),
the dot is returned to the input and the integer token is returned, so
\lain{a..b} lexes correctly as an integer followed by the range operator.

\subsection{String literals}
\label{sec:lex-strings}
\index{string literals!lexing}

\pnum
In \cname{STATE\_DOUBLE\_QUOTE}, a backslash causes the FSM to advance
\cname{current} one extra byte, consuming the escape without interpreting it.
Escape decoding (\texttt{\\n}, \texttt{\\t}, \texttt{\\0}, etc.) is performed
later in the parser.  The resulting token's \cname{start} points past the opening
quote and \cname{length} excludes the closing quote.

\subsection{Nested block comments}
\label{sec:lex-comments}
\index{comments!nested block}

\pnum
Block comments (\texttt{/* ... */}) are handled inline in the
\cname{STATE\_SLASH} case.  A \texttt{nesting} counter supports arbitrarily deep
nesting:

\begin{ccode}
int nesting = 1;
while (nesting > 0) {
    if (d == '/' && current[1] == '*')       { nesting++; current += 2; }
    else if (d == '*' && current[1] == '/')  { nesting--; current += 2; }
    else if (d == 0) break;  /* EOF */
    else current++;
}
\end{ccode}

% -------------------------------------------------------------------------
\section{Keyword classification}
\label{sec:lex-keywords}
\index{keywords!classification}
% -------------------------------------------------------------------------

\pnum
After an identifier token is scanned, \cname{token\_match\_keyword(lexeme, len)}
(\srcref{src/token.h:103}) tests whether it is a reserved keyword.  The function
dispatches on length, then uses \cname{strncmp}:

\begin{center}
\begin{tabular}{rl}
\toprule
\textbf{Len} & \textbf{Keywords} \\
\midrule
2  & \lain{if}, \lain{in}, \lain{as}, \lain{or} \\
3  & \lain{end}, \lain{for}, \lain{var}, \lain{mov}, \lain{use}, \lain{and}, \lain{pre} \\
4  & \lain{type}, \lain{func}, \lain{proc}, \lain{expr}, \lain{elif}, \lain{else},
     \lain{case}, \lain{post}, \lain{true} \\
5  & \lain{break}, \lain{false}, \lain{macro}, \lain{while}, \lain{defer} \\
6  & \lain{import}, \lain{export}, \lain{extern}, \lain{return}, \lain{unsafe} \\
8  & \lain{continue}, \lain{comptime} \\
9  & \lain{c\_include}, \lain{undefined} \\
10 & \lain{decreasing} \\
\bottomrule
\end{tabular}
\end{center}

\pnum
If no keyword matches, \cname{TOKEN\_IDENTIFIER} is returned.  The length
dispatch ensures that most identifiers avoid scanning the entire keyword table.

% -------------------------------------------------------------------------
\section{Lookahead}
\label{sec:lex-peek}
\index{lexer!lookahead}
% -------------------------------------------------------------------------

\pnum
\cname{lexer\_peek(Lexer *lexer)} (\srcref{src/lexer.h:380}) provides
non-consuming lookahead by copying the lexer state, calling \cname{lexer\_next()}
on the copy, and returning the result without modifying the original:

\begin{ccode}
Token lexer_peek(Lexer *lexer) {
    Lexer temp = *lexer;
    return lexer_next(&temp);
}
\end{ccode}

Because \cname{Lexer} contains only two pointers, the copy is eight bytes on a
64-bit system and the peek has the same cost as a regular \cname{lexer\_next()}.
